\chapter{The Exponential and the Logarithm}\label{ch:explog}

$e^{\ln x} = x$ is the definition of the logarithm, and the engine's proof of it fails at
$x = -1$. The reason is that Mathlib's logarithm is a total function, defined for negative numbers
as $\ln |x|$, so that $\ln(-1) = 0$ and $e^{\ln(-1)} = 1$. This chapter relearns the two functions
from their definitions, states their identities with their domains, and shows which of the
engine's function rules are proved, which is refuted, and which have no theorem yet.

\section{The exponential}

\begin{definition}
  $\exp(x) = \sum_{n \ge 0} \frac{x^n}{n!}$.
\end{definition}

The series converges absolutely for every real (and complex) $x$ by the ratio test:
$\frac{|x|^{n+1}/(n+1)!}{|x|^n/n!} = \frac{|x|}{n+1} \to 0$.

\begin{theorem}
  $\exp(x + y) = \exp(x)\exp(y)$; $\exp(0) = 1$; $\exp$ is differentiable with $\exp' = \exp$;
  $\exp(x) > 0$ for all real $x$; $\exp$ is strictly increasing and maps $\mathbb{R}$ onto
  $(0, \infty)$.
\end{theorem}
\begin{proof}[Proof sketch]
  The functional equation is the Cauchy product of two absolutely convergent series with the
  binomial theorem (Chapter~\ref{ch:poly}) applied termwise: the coefficient of $\frac{1}{n!}$ in
  the product is $\sum_k \binom{n}{k} x^k y^{n-k} = (x+y)^n$. Differentiating termwise, which is
  legitimate for a power series inside its radius of convergence, gives $\exp' = \exp$.
  Positivity: $\exp(x)\exp(-x) = \exp(0) = 1$, so $\exp(x) \neq 0$, and $\exp(x) > 0$ for
  $x \ge 0$ from the series; then for $x < 0$ from $\exp(x) = 1/\exp(-x)$. Strictly increasing
  since the derivative is positive; onto $(0, \infty)$ by the intermediate value theorem with
  $\exp(x) \ge 1 + x$ and $\exp(-x) \le 1/(1+x)$.
\end{proof}

\section{The logarithm}

\begin{definition}
  $\ln : (0, \infty) \to \mathbb{R}$ is the inverse of $\exp$: $\ln y$ is the unique $x$ with
  $\exp x = y$.
\end{definition}

\begin{theorem}
  For $a, b > 0$: $\ln(ab) = \ln a + \ln b$; $\ln 1 = 0$; $\ln(a^p) = p \ln a$ for real $p$;
  $\ln$ is differentiable with $\ln'(y) = 1/y$.
\end{theorem}
\begin{proof}
  $\exp(\ln a + \ln b) = \exp(\ln a)\exp(\ln b) = ab$, so $\ln a + \ln b$ is the $x$ with
  $\exp x = ab$. $\ln 1 = 0$ since $\exp 0 = 1$. With $a^p := \exp(p \ln a)$, $\ln(a^p) = p \ln a$
  is immediate from the definition. The derivative is the inverse function rule:
  $\ln'(y) = 1/\exp'(\ln y) = 1/\exp(\ln y) = 1/y$.
\end{proof}

Every identity has the hypothesis $a, b > 0$, because that is the domain of $\ln$.

\section{Mathlib's logarithm}

Mathlib defines \lc{Real.log x} for every real $x$: it is $\ln |x|$ for $x \neq 0$, and $0$ at
$0$. The consequences are three lemmas the engine's proofs use.

\begin{itemize}[nosep]
  \item \lc{Real.log_exp} : $\ln(\exp x) = x$ for all $x$. Unconditional, because $\exp x$ is
    always positive.
  \item \lc{Real.exp_log} : $\exp(\ln x) = x$ for $0 < x$. Conditional. For $x < 0$,
    $\exp(\ln x) = \exp(\ln|x|) = |x| = -x$.
  \item \lc{Real.log_neg_eq_log} : $\ln(-x) = \ln x$. The logarithm is an even function.
\end{itemize}

The third is the one that makes the counterexample two lines long.

\section{The rules}

The engine's function rules are a table of identities at special arguments and inverse pairs:

\begin{minted}{lean}
def functionApply : Expr → Option RuleResult
  | .mul es =>
    match findTan es with
    | some (u, others) => some ⟨mulN (.fn "tan" [u] :: others), "$\\sin u / \\cos u = \\tan u$.", none, none⟩
    | none => none
  | .fn "sqrt" [a] => some ⟨.pow a (.num (Q.ofRat (mkRat 1 2))), "…", none, none⟩
  | .fn "ln" [a] =>
    if isOne a then some ⟨Expr.zero, "$\\ln 1 = 0$.", none, none⟩ else
    match a with
    | .fn "exp" [x] => some ⟨x, "$\\ln(e^x) = x$: ln and exp are inverses.", none, none⟩
    | .pow b p => some ⟨.mul [p, .fn "ln" [b]], "$\\ln(b^p) = p \\ln b$.", none, none⟩
    | _ => none
  | .fn "exp" [a] =>
    if isZero a then some ⟨Expr.one, "$e^0 = 1$.", none, none⟩ else
    match a with
    | .fn "ln" [x] => some ⟨x, "$e^{\\ln x} = x$: exp and ln are inverses.", none, none⟩
    | _ => none
  | .fn "sin" [a] => if isZero a then some ⟨Expr.zero, "$\\sin 0 = 0$.", none, none⟩ else none
  | .fn "cos" [a] => if isZero a then some ⟨Expr.one, "$\\cos 0 = 1$.", none, none⟩ else none
  | .fn "abs" [.num q] => some ⟨.num q.abs, "Absolute value of a constant.", none, none⟩
  | _ => none
\end{minted}

\section{The theorems, and the gaps}

Two cases are settled by theorems in \file{proofs/}.

\begin{thmbox}[\lcn{not_functionRules_soundR}]
  The rule is not unconditionally sound: at $x = -1$ it rewrites $\exp(\ln x)$, whose value is
  $\exp(\ln 1) = 1$, to $x$, whose value is $-1$.
\end{thmbox}

\begin{minted}{lean}
theorem not_functionRules_soundR : ¬ RuleSoundR functionRules := by
  intro hs
  have h := hs (.fn "exp" [.fn "ln" [.var "x"]]) _ rfl (fun _ => -1)
  simp only [evalR_fn₁, applyFn_exp, applyFn_ln, evalR_var] at h
  rw [show ((-1 : ℝ)) = -(1 : ℝ) by norm_num, Real.log_neg_eq_log, Real.log_one, Real.exp_zero] at h
  norm_num at h
\end{minted}

The positive theorem for that case is \lc{exp_log_sound}, which is \lc{Real.exp_log} with the
hypothesis $0 < x$; the notebook's label for the rule is \status{conditional} with that
hypothesis.

\begin{thmbox}[\lcn{functionRules_tan_soundR}]
  The case $\sin u \cdot (\cos u)^{-1} \to \tan u$ is unconditionally sound.
\end{thmbox}

This is the case that was added when the integral checker refused $\int \tan x$
(Chapter~\ref{ch:parts}). It is unconditional because Mathlib \emph{defines} $\tan x$ as
$\sin x / \cos x$ with $x/0 = 0$, so the two sides agree even where the cosine vanishes; the
proof is \lc{Real.tan_eq_sin_div_cos} and the permutation lemma for products, since the two
factors may sit anywhere in the list.

The remaining cases have no theorem. $\ln 1 = 0$, $e^0 = 1$, $\sin 0 = 0$, $\cos 0 = 1$,
$\ln(e^x) = x$, $\sqrt a = a^{1/2}$ and $|c|$ for a numeral are all true unconditionally in
Mathlib's semantics (\lc{Real.log_one}, \lc{Real.exp_zero}, \lc{Real.sin_zero},
\lc{Real.cos_zero}, \lc{Real.log_exp}, \lc{Real.sqrt_eq_rpow}, \lc{Rat.cast_abs}) and would each
be a one-line theorem. $\ln(b^p) \to p \ln b$ is \emph{not} unconditional: for $b = -2$ and
$p = 1/2$, Mathlib's $(-2)^{1/2} = \sqrt 2 \cos(\pi/2) = 0$, so the left side is $\ln 0 = 0$ while
the right side is $\frac{1}{2}\ln 2$. It needs $0 < b$ (\lc{Real.log_rpow}), and it needs a
counterexample theorem like the one above. The rule table's note, ``the other cases are
unconditional'', is therefore an overstatement for that one case, and the ledger should say
\status{conditional} for it once the theorems are written. This is recorded here rather than
hidden: the book's claim is not that everything is proved, but that the boundary is visible.

\begin{trybox}
\begin{minted}{text}
exp(ln(x))
ln(exp(x))
sin(x) / cos(x)
ln(x^2)
\end{minted}
  The first is labelled with its condition. The second has no condition, and no theorem yet.
  The third is the tangent case. The fourth fires the one case that should be conditional and
  is not yet labelled so.
\end{trybox}

\section{Exercises}

\begin{exercisebox}[The gaps]
  Write and prove \lc{functionRules_ln_exp_soundR} in the style of
  \lc{functionRules_tan_soundR}. Then state the counterexample theorem for $\ln(b^p) \to
  p \ln b$ and prove it with \lc{Real.rpow_def_of_neg}.
\end{exercisebox}

\begin{exercisebox}[The functional equation]
  Fill in the Cauchy-product step of the proof that $\exp(x+y) = \exp(x)\exp(y)$: write the
  product of the two series as a double sum, regroup by $n = j + k$, and identify the binomial
  theorem.
\end{exercisebox}

\begin{exercisebox}[Evenness]
  Explain why defining $\ln x = \ln |x|$ for $x < 0$ is the natural choice for a total function
  in the sense that $\ln$ remains a local inverse of $\exp$ on $\exp$'s image and satisfies
  $\ln(ab) = \ln a + \ln b$ for all nonzero $a, b$. Find the identity of the theorem above that
  it does \emph{not} preserve.
\end{exercisebox}
